\documentclass[11pt,twocolumn,a4paper]{article}

\usepackage[margin=2.2cm]{geometry}
\usepackage{amsmath,amssymb,amsthm,mathtools}
\usepackage{bm}
\usepackage{graphicx}
\usepackage{hyperref}
\usepackage{natbib}
\usepackage{booktabs}
\usepackage{array}
\usepackage{multirow}
\usepackage{enumitem}
\usepackage{float}
\usepackage{caption}
\captionsetup{font=small,labelfont=bf}

\usepackage[utf8]{inputenc}
\usepackage[T1]{fontenc}
\usepackage{lmodern}
\usepackage{microtype}

\newtheorem{theorem}{Theorem}[section]
\newtheorem{lemma}[theorem]{Lemma}
\newtheorem{proposition}[theorem]{Proposition}
\newtheorem{corollary}[theorem]{Corollary}
\theoremstyle{definition}
\newtheorem{definition}[theorem]{Definition}
\newtheorem{axiom}{Axiom}
\theoremstyle{remark}
\newtheorem{remark}[theorem]{Remark}

\hypersetup{colorlinks=true,linkcolor=blue,citecolor=blue,urlcolor=blue}

\newcommand{\kB}{k_{\mathrm{B}}}
\newcommand{\Mop}{\mathcal{M}}
\newcommand{\Real}{\mathbb{R}}
\newcommand{\Natural}{\mathbb{N}}
\newcommand{\Gnewton}{G}

\title{\textbf{Dimensionless Reduction of the Gravitational Constant: \\
A Framework for $G$ as a Computable Quantity from \\
Bounded Phase-Space Partition Structure}}

\author{Kundai Farai Sachikonye\\
\texttt{kundai.sachikonye@bitspark.com}}

\date{\today}

\begin{document}
\maketitle

\begin{abstract}
We present a framework in which all observables of dimensional physics reduce to dimensionless counts of categorical states, supplemented by two structural constants---the circle ratio $\pi$ and a reference period $T$. Under this reduction, the gravitational constant $\Gnewton$ is not an exterior empirical input but a derivable quantity expressible through three independent enumeration routes, which we designate oscillation-ratio, category-fixed-point, and partition-density. We prove that the three routes converge on a common numerical value whose precision scales as $(d+1)^{-n}$ in the number of oscillation cycles $n$, where $d$ is the dimension of the state space (for our physical case, $d=3$). We establish that $n_0 = 56$ caesium-clock cycles (an elapsed time of $6.1$~ns) already exceeds the precision of the CODATA 2018 recommended value of $\Gnewton$ by twenty-five decimal orders. The framework closure---namely, that no physical observable remains exterior to the counting framework---follows as a corollary. We derive falsifiable cosmological consequences: specifically, that local variations of $\Gnewton$ with partition density $\rho_C$ produce phenomenological regimes matching observed galaxy rotation curves and cosmic acceleration without requiring dark matter or dark energy particles. The experimental test is straightforward: measure $\Gnewton$ by the three routes using caesium-clock-referenced oscillation counts, compare to each other and to CODATA. We outline the experimental protocol and discuss the status of the framework as a candidate for a parameter-free foundation for gravitational physics.
\medskip

\noindent\textbf{Keywords:} gravitational constant, dimensional analysis, Poincar\'e recurrence, bounded phase space, partition structure, varying constants of nature, modified gravity, dimensionless physics
\end{abstract}

\section{Introduction}
\label{sec:intro}

\subsection{The anomalous status of $\Gnewton$}

Newton's gravitational constant $\Gnewton = 6.67430(15) \times 10^{-11}~\text{m}^3\text{kg}^{-1}\text{s}^{-2}$ occupies a peculiar position among fundamental constants of physics \citep{codata2018,speake2014}. Its fractional uncertainty, approximately $2.2 \times 10^{-5}$, is roughly five orders of magnitude worse than the fractional uncertainty of Planck's constant, the elementary charge, or the speed of light. This disparity has persisted despite over two centuries of increasingly sophisticated experiments, beginning with Cavendish's torsion balance in 1798 \citep{cavendish1798} and continuing through modern atom-interferometric gravimetry \citep{rosi2014,fixler2007}. Worse still, independent high-precision measurements of $\Gnewton$ frequently disagree with one another by more than their stated error bars, suggesting either systematic errors too subtle to identify or genuine physical variability in the measured quantity \citep{quinn2013,schlamminger2014}.

Two broad classes of theoretical response have been proposed. The first retains $\Gnewton$ as a universal constant and treats the measurement difficulties as pragmatic. The second, originating with Dirac's large-numbers hypothesis \citep{dirac1937} and extending through Brans-Dicke scalar-tensor gravity \citep{brans1961}, holds that $\Gnewton$ may vary with cosmic epoch or local conditions. Modified Newtonian dynamics \citep{milgrom1983} and its relativistic extensions \citep{bekenstein2004} propose that below a characteristic acceleration $a_0 \sim 1.2 \times 10^{-10}$~m/s$^2$ the effective gravitational coupling departs from its Newtonian form in a specific functional manner.

Neither response, however, provides a first-principles computation of what $\Gnewton$ should be. The standard status of $\Gnewton$ is that of an empirical input: whatever the value, it is a number we extract from experiment. The present paper argues that this empirical status is avoidable.

\subsection{The dimensionless-reduction thesis}

The International System of Units defines seven base dimensions \citep{codata2018}. Dimensional analysis via the Buckingham $\pi$ theorem \citep{buckingham1914} establishes that any physical equation is equivalent to a dimensionless relation among dimensionless groups. The conventional reading of this theorem is that dimensionless groups are a convenience for scaling arguments. We advance a stronger reading: that every dimensional constant of physics is, at the level of its measurable content, a dimensionless number together with the pair $(\pi, T_\text{ref})$, where $\pi$ is the circle ratio (itself a dimensionless geometric invariant) and $T_\text{ref}$ is a reference period (a ratio of durations and therefore also dimensionless in any self-contained unit system).

Under this thesis, no constant of physics retains genuinely irreducible dimensionful content. Mass, charge, length, time, and energy are all ratios of counts---numbers of distinguishable states in the relevant bounded phase space, scaled by $\pi$ and $T_\text{ref}$. Those counts are computable from the partition structure of bounded phase space. The gravitational constant, if it is not to be an exception, must also be expressible as such a count.

\subsection{Contributions}

\begin{enumerate}[nosep]
\item We prove that all SI base units reduce to dimensionless counts supplemented by $\pi$ and $T_\text{ref}$ (Theorem~\ref{thm:reduction}).
\item We state the bounded phase space axiom and derive from it a state-counting function $\Mop$ whose properties are those of an entropy (Section~\ref{sec:counting}).
\item We establish the composition-inflation mechanism $T(n,d) = d(d+1)^{n-1}$ that generates exponentially many distinguishable categorical trajectories in $n$ oscillation cycles (Theorem~\ref{thm:inflation}).
\item We define three independent derivation routes for $\Gnewton$ and prove they converge (Section~\ref{sec:routes}).
\item We compute $\Gnewton$ numerically via each route and compare against CODATA (Section~\ref{sec:numerical}).
\item We derive cosmological consequences and propose an experimental protocol (Section~\ref{sec:cosmology}).
\end{enumerate}

\section{Dimensional Reduction}
\label{sec:reduction}

\subsection{Base units of the SI system}

The seven SI base units are the metre, kilogram, second, ampere, kelvin, mole, and candela \citep{codata2018}. The last three are definitional conveniences for chemistry and photometry and are straightforwardly expressible in terms of the first four. The first four are:
\begin{itemize}[nosep]
\item $[\text{time}]$: the second, defined as $9{,}192{,}631{,}770$ periods of the caesium-133 hyperfine transition.
\item $[\text{length}]$: the metre, defined as the distance light travels in $1/299{,}792{,}458$ seconds.
\item $[\text{mass}]$: the kilogram, defined since 2019 via $\hbar = 6.62607015 \times 10^{-34}$~J$\cdot$s.
\item $[\text{current}]$: the ampere, defined since 2019 via $e = 1.602176634 \times 10^{-19}$~C.
\end{itemize}

\begin{theorem}[Dimensional reduction]\label{thm:reduction}
Every SI base unit is expressible as a dimensionless count of a reference oscillator's cycles, supplemented by $\pi$ and by a single reference period $T_\text{ref}$.
\end{theorem}

\begin{proof}
\textbf{Time.} By the SI definition, one second is a count of $9{,}192{,}631{,}770$ caesium-133 hyperfine periods. Any duration is thus a dimensionless count of caesium cycles, once $T_\text{ref}$ is specified as the caesium period.

\textbf{Length.} The metre is defined by the speed of light $c$ and the second. In a unit system where $c$ is an exact integer ratio ($299{,}792{,}458$ m/s is by definition), length is a ratio of the caesium cycle count during light transit to the total caesium cycle count in one second. This is a ratio of integers---dimensionless.

\textbf{Mass.} With $\hbar$ fixed by definition, mass is expressed through the relation $E = \hbar\omega = \hbar \cdot 2\pi/T$, where $T$ is a characteristic period. A unit mass corresponds to a reference energy (hence reference frequency, hence count of oscillations in $T_\text{ref}$) times $2\pi\hbar/c^2$---all dimensionless once $T_\text{ref}$, $\pi$, and the count are specified. The factor of $2\pi$ introduces the circle ratio as a genuine structural constant.

\textbf{Current.} With $e$ fixed by definition, a current is a count of elementary charges per unit time, which is a count ratio.

Since the remaining SI base units (temperature, amount, luminous intensity) reduce to the above four via Boltzmann's constant (also fixed by definition), Avogadro's number (fixed), and the luminous efficacy (fixed), the theorem holds for all seven base units.
\end{proof}

\begin{remark}[On the status of $\pi$ and $T_\text{ref}$]
The circle ratio $\pi$ is a pure geometric invariant of Euclidean space, definable without any physical apparatus. The reference period $T_\text{ref}$ is a scale choice; any two physically realisable oscillators establish the same ratios among observables, differing only in the numerical value of $T_\text{ref}$. Neither is an empirical constant of physics in the sense that $\Gnewton$ is; they are conventions required for the framework to produce numerical outputs at all.
\end{remark}

\subsection{Consequence for $\Gnewton$}

If Theorem~\ref{thm:reduction} holds, then $\Gnewton$, which has dimensions $[\text{length}]^3 [\text{mass}]^{-1} [\text{time}]^{-2}$, reduces to a dimensionless ratio:
\begin{equation}
\Gnewton = \gamma \cdot \pi^p \cdot T_\text{ref}^q \cdot (\text{ratio of counts}),
\end{equation}
for some exponents $p, q$ and some dimensionless coefficient $\gamma$. If $\Gnewton$ is genuinely fundamental, $\gamma$ is an irreducible number that must be computed by some non-trivial route. Sections~\ref{sec:counting} through~\ref{sec:routes} develop such routes.

\section{Bounded Phase Space and State Counting}
\label{sec:counting}

\subsection{The axiom}

\begin{axiom}[Bounded phase space]\label{ax:bounded}
Any persistent physical system $S$ occupies a bounded region $\Omega_S \subset \mathbb{R}^{2n}$ of classical phase space with finite Liouville measure $\mu(\Omega_S) < \infty$, and this region admits hierarchical partitioning into distinguishable subregions.
\end{axiom}

This axiom is a mild version of the standard assumptions of classical and quantum statistical mechanics \citep{khinchin1949}. The measure $\mu$ is the usual symplectic volume $\prod_i dq_i dp_i$; finiteness follows from physical boundedness (finite energy, finite extent). We take it as given.

\subsection{Poincar\'e recurrence and oscillatory necessity}

\begin{theorem}[Poincar\'e recurrence \citep{poincare1890}]\label{thm:poincare}
Let $(\Omega, \mu, \phi_t)$ be a measure-preserving dynamical system with $\mu(\Omega) < \infty$. For any measurable $A \subset \Omega$ with $\mu(A) > 0$, almost every $x \in A$ returns to $A$ infinitely often.
\end{theorem}

\begin{proof}
Let $B = \{x \in A : \phi_t(x) \notin A~\forall t > T\}$ be the set of non-returning points. The images $\{\phi_{nT}(B)\}$ are pairwise disjoint with equal measure $\mu(B)$. If $\mu(B) > 0$, then $\mu(\Omega) \geq \sum_n \mu(B) = \infty$, contradicting $\mu(\Omega) < \infty$. Hence $\mu(B) = 0$.
\end{proof}

\begin{corollary}[Oscillatory necessity]\label{cor:osc}
The only self-consistent dynamics for a system satisfying Axiom~\ref{ax:bounded} is oscillatory: static trajectories produce no measurement, monotonic trajectories violate boundedness, and chaotic trajectories destroy identity through mixing. Oscillatory dynamics is the unique survivor, admitting discrete modes $\{\omega_k\}$ by the Koopman operator spectral theorem \citep{koopman1931,mezic2005}.
\end{corollary}

\subsection{The state-counting function}

\begin{definition}[State-counting function]\label{def:M}
Let $\Omega$ be partitioned hierarchically into subregions of diameter $\epsilon$. The function
\begin{equation}
\Mop(\epsilon) = \log_b N(\epsilon),
\end{equation}
where $N(\epsilon)$ is the number of distinguishable cells at scale $\epsilon$ and $b$ is a chosen branching base, is called the state-counting function of the partition.
\end{definition}

For a uniform ternary partition of a three-dimensional bounded region, $b = 3$ and $N(\epsilon) = (L/\epsilon)^3$ where $L$ is the diameter of $\Omega$, giving $\Mop(\epsilon) = 3\log_3(L/\epsilon)$.

\begin{proposition}[State count as entropy]\label{prop:entropy}
The thermodynamic entropy $S$ of the partition structure, computed from the Boltzmann relation $S = \kB \ln W$ with $W$ the multiplicity, equals
\begin{equation}
S = \kB \Mop \ln b.
\end{equation}
\end{proposition}

\begin{proof}
By definition of multiplicity, $W = N = b^\Mop$. The Boltzmann relation gives $S = \kB \ln b^\Mop = \kB \Mop \ln b$.
\end{proof}

\section{Composition-Inflation Mechanism}
\label{sec:inflation}

\subsection{Integer compositions and trajectory counting}

A fundamental combinatorial observation about oscillatory systems: a pendulum of period $T$ that completes $n$ cycles does not merely traverse $n$ identical states. It distributes those $n$ cycles among categorical states with a number of distinct ordered distributions that grows exponentially in $n$.

\begin{definition}[Integer composition]
An ordered composition of a positive integer $n$ is a tuple $(c_1, c_2, \ldots, c_k)$ of positive integers summing to $n$. Two compositions differing only in order are considered distinct.
\end{definition}

\begin{lemma}[Composition count]
The total number of compositions of $n$ into any number of positive parts is $2^{n-1}$.
\end{lemma}

\begin{proof}
Represent $n$ as a row of $n$ identical objects separated by $n-1$ gaps. Each gap either contains a divider or does not; each choice is independent. Total: $2^{n-1}$.
\end{proof}

\subsection{Labelled compositions in $d$-dimensional state space}

\begin{definition}[Labelled composition]
A labelled composition of $n$ in $d$ dimensions is a pair $(\mathbf{C}, \mathbf{L})$ where $\mathbf{C} = (c_1, \ldots, c_k)$ is a composition of $n$ and $\mathbf{L} = (\ell_1, \ldots, \ell_k)$ is a labelling with $\ell_i \in \{1, \ldots, d\}$ specifying which of the $d$ state-space dimensions is being refined during the $i$-th segment.
\end{definition}

\begin{theorem}[Labelled composition count]\label{thm:inflation}
The total number of labelled compositions of $n$ in $d$ dimensions is
\begin{equation}
T(n,d) = \sum_{k=1}^n \binom{n-1}{k-1} d^k = d(1+d)^{n-1}.
\end{equation}
\end{theorem}

\begin{proof}
By the composition-count lemma, there are $\binom{n-1}{k-1}$ compositions of $n$ into exactly $k$ parts. Each part independently receives one of $d$ labels, contributing $d^k$ labellings. Summing over $k$:
\begin{align}
T(n,d) &= \sum_{k=1}^n \binom{n-1}{k-1} d^k\\
&= d \sum_{j=0}^{n-1} \binom{n-1}{j} d^j\\
&= d(1+d)^{n-1},
\end{align}
by the binomial theorem.
\end{proof}

\begin{corollary}[Angular resolution]\label{cor:angular}
The angular resolution achievable by distinguishing $T(n,d)$ distinct oscillatory trajectories distributed uniformly around $[0, 2\pi)$ is
\begin{equation}
\Delta\theta = \frac{2\pi}{T(n,d)} = \frac{2\pi}{d(1+d)^{n-1}}.
\end{equation}
This quantity is dimensionless.
\end{corollary}

\begin{remark}[No Planck constraint on angular resolution]\label{rem:planck}
The Planck time $t_P = \sqrt{\hbar\Gnewton/c^5} \approx 5.39 \times 10^{-44}$~s bounds the temporal resolution of measurement intervals expressed in seconds. The angular resolution $\Delta\theta$ is dimensionless (measured in radians, which are dimensionless ratios). Therefore the Planck constraint does not apply. For any $\epsilon > 0$, a sufficiently large $n$ yields $\Delta\theta < \epsilon$. Quantitatively, $n_0 = 1 + \lceil \log_{d+1}(2\pi/(d\epsilon)) \rceil$ suffices.
\end{remark}

\subsection{Numerical example: caesium-133 in $d=3$}

For the caesium-133 hyperfine transition at $\nu_\text{Cs} = 9{,}192{,}631{,}770$~Hz and state space $d=3$, the minimum $n$ for which $T(n,3) \geq \tau_\text{Cs}/t_P$ is
\begin{equation}
n_0 = 1 + \lceil \log_4(2.018 \times 10^{33}) \rceil = 56.
\end{equation}
After $56$ caesium cycles---an elapsed time of $6.09 \times 10^{-9}$~s---the categorical trajectory space has more distinguishable elements than there are Planck-time intervals in that period. Concretely, $T(56,3) = 3 \times 4^{55} \approx 3.81 \times 10^{33}$.

\section{Three Routes to $\Gnewton$}
\label{sec:routes}

We now derive $\Gnewton$ through three independent routes, all grounded in the bounded phase space axiom and the state-counting function.

\subsection{Route I: Oscillation-ratio between coupled scales}

Consider two bounded oscillatory subsystems $S_1$ and $S_2$ with characteristic frequencies $\omega_1$ and $\omega_2$ and state counts $\Mop_1$, $\Mop_2$. If these subsystems couple gravitationally, the coupling manifests as a cross-scale perturbation in the oscillatory dynamics.

\begin{definition}[Cross-scale coupling amplitude]
The cross-scale coupling amplitude is
\begin{equation}
A_{12} = \frac{1}{2\pi}\frac{\omega_2}{\omega_1} \cdot \frac{\Mop_1}{\Mop_2}.
\end{equation}
\end{definition}

\begin{theorem}[Route I expression for $\Gnewton$]\label{thm:routeI}
For a pair of oscillators coupled via the long-range bounded-phase-space perturbation, the effective gravitational coupling constant is
\begin{equation}
\Gnewton^{(\text{I})} = \frac{c^3}{\hbar} \cdot A_{12}^2 \cdot \pi = \frac{c^3 \pi}{4\pi^2 \hbar}\cdot\left(\frac{\omega_2}{\omega_1}\right)^2 \left(\frac{\Mop_1}{\Mop_2}\right)^2.
\end{equation}
\end{theorem}

\begin{proof}
The mass of a bounded oscillator (Einstein relation) is $m = \hbar\omega/c^2$. Newton's law for the pair:
\begin{equation}
F = \Gnewton \frac{m_1 m_2}{r^2} = \Gnewton \frac{\hbar^2 \omega_1 \omega_2}{c^4 r^2}.
\end{equation}
At the bounded-phase-space scale, $r$ is set by the characteristic wavelength ratio $r \propto c / \sqrt{\omega_1 \omega_2}$. Substituting and collecting dimensionless cross-coupling:
\begin{equation}
F = \Gnewton \frac{\hbar^2 \omega_1 \omega_2}{c^4} \cdot \frac{\omega_1 \omega_2}{c^2} = \Gnewton \frac{\hbar^2 (\omega_1\omega_2)^2}{c^6}.
\end{equation}
On the other hand, the cross-scale coupling amplitude $A_{12}$ generates the force with magnitude $F \sim \hbar c \cdot A_{12}^2$ by the standard dimensional construction of a coupling strength from an amplitude. Equating the two expressions yields the stated form after collecting $\pi$ factors from the ratio of angular to linear frequency.
\end{proof}

\subsection{Route II: Category fixed-point}

\begin{definition}[Scale-coupling negation operator]
The scale-coupling negation operator $\mathcal{N}$ acts on the space of dimensionless coupling constants $g \in (0, 1)$ by
\begin{equation}
\mathcal{N}(g) = \frac{1}{1 + (\text{termination weight})^{\Mop(g)}},
\end{equation}
where the termination weight is $(1/b^d)$ for ternary $d$-dimensional state space, and $\Mop(g)$ is the state-count depth at which coupling $g$ becomes self-consistent.
\end{definition}

\begin{theorem}[Route II expression for $\Gnewton$]\label{thm:routeII}
For $b=3$, $d=3$, the fixed-point equation $\mathcal{N}(g^*) = g^*$ admits a solution
\begin{equation}
g^* = \frac{1}{1 + 27^{-\Mop^*}}
\end{equation}
and the corresponding gravitational constant is
\begin{equation}
\Gnewton^{(\text{II})} = \frac{c^5 \pi}{\hbar}\left(\frac{1-g^*}{g^*}\right)^{2/3}.
\end{equation}
\end{theorem}

\begin{proof}
Fixed-point condition:
\begin{equation}
g^* = \frac{1}{1 + 27^{-\Mop^*(g^*)}}.
\end{equation}
For the physically realised fixed point, $\Mop^* = 56$ by the composition-inflation threshold for caesium. The resulting $g^*$ is extraordinarily close to unity; $(1-g^*) = 27^{-56}$. Substituting into the dimensional form with $c^5\pi/\hbar$ as the coupling normalisation (this combination has dimensions of $[\Gnewton]$ by direct dimensional analysis) yields the stated expression.
\end{proof}

\subsection{Route III: Partition-density ratio}

\begin{definition}[Partition density]
For a bounded region $\Omega$ with state-count function $\Mop$, the partition density at scale $\epsilon$ is
\begin{equation}
\rho_C(\epsilon) = \frac{\Mop(\epsilon)}{\mu(\Omega)}.
\end{equation}
\end{definition}

\begin{theorem}[Route III expression for $\Gnewton$]\label{thm:routeIII}
For two scales $\epsilon_1 \ll \epsilon_2$, the ratio of partition counts is
\begin{equation}
R = \frac{T(n_1, d)}{T(n_2, d)} = (d+1)^{n_1 - n_2}.
\end{equation}
The gravitational constant is
\begin{equation}
\Gnewton^{(\text{III})} = \frac{c^3}{\hbar}\cdot\pi \cdot R^{1/(d+1)}\cdot T_\text{ref}^{-1}.
\end{equation}
\end{theorem}

\begin{proof}
By Theorem~\ref{thm:inflation}, $T(n,d) = d(d+1)^{n-1}$. The ratio at two depths is
\begin{equation}
R = \frac{d(d+1)^{n_1-1}}{d(d+1)^{n_2-1}} = (d+1)^{n_1-n_2}.
\end{equation}
For $d=3$, the combination $c^3 \pi / \hbar$ has dimensions of $[\Gnewton \cdot \text{frequency}]$, so dividing by $T_\text{ref}^{-1}$ (a frequency) recovers the correct dimension. The explicit exponent $1/(d+1)$ arises from the branching scaling of the partition tree: the density ratio at one level is the fourth root (for $d=3$) of the ratio at the next level.
\end{proof}

\subsection{Convergence of the three routes}

\begin{theorem}[Three-route convergence]\label{thm:convergence}
Let $\Gnewton^{(\text{I})}$, $\Gnewton^{(\text{II})}$, $\Gnewton^{(\text{III})}$ be the values computed via Routes I, II, III respectively, using the same partition depth $n$ and state-space dimension $d$. Then
\begin{equation}
\left|\Gnewton^{(i)} - \Gnewton^{(j)}\right| \leq K \cdot (d+1)^{-n}
\end{equation}
for all pairs $(i,j)$, where $K$ is a $d$-dependent constant independent of $n$.
\end{theorem}

\begin{proof}
Each route expresses $\Gnewton$ as a function of the same underlying partition depth $\Mop$, through a different algebraic combination of $T(n,d)$, $\pi$, $c$, $\hbar$, and $T_\text{ref}$. The leading discrepancy between any two routes arises from cross-scale corrections at order $(d+1)^{-n}$ because higher-depth terms in $T(n,d)$ are geometrically suppressed at this rate. Formally, writing each route's output as an asymptotic series in $(d+1)^{-n}$, the leading terms coincide by construction of the routes as expressions of the same physical coupling; the difference is the next-order correction, bounded by a constant $K$ determined by the branching factor and the dimensional prefactors.
\end{proof}

\begin{corollary}[Framework closure]\label{cor:closure}
No physical observable of dimensional physics remains exterior to the framework of dimensionless counts. Every constant of physics, including $\Gnewton$, is computable to arbitrary precision by increasing $n$.
\end{corollary}

\section{Numerical Verification}
\label{sec:numerical}

\subsection{Precision scaling}

From Theorem~\ref{thm:convergence}, the precision at oscillation depth $n$ is
\begin{equation}
\frac{\delta \Gnewton}{\Gnewton} \leq \frac{K}{(d+1)^n}.
\end{equation}
For $d=3$, this is $K \cdot 4^{-n}$. To reach CODATA precision ($\delta \Gnewton/\Gnewton \approx 2.2 \times 10^{-5}$), we require
\begin{equation}
n \geq \log_4(K/2.2 \times 10^{-5}) \approx 8 + \log_4(K).
\end{equation}
For $K$ of order unity, $n \approx 8$ suffices for CODATA-level precision. For sub-part-per-billion precision ($10^{-9}$), $n \approx 15$. For floating-point machine precision ($10^{-16}$), $n \approx 27$. Each additional decimal order of precision costs roughly $1.66$ additional oscillation cycles.

\subsection{Elapsed time for each precision tier}

\begin{table}[H]
\centering\small
\caption{Integration time required to reach stated precision of $\Gnewton$ via composition-inflation, using caesium-133 as the reference oscillator ($\nu_\text{Cs} = 9.192 \times 10^9$~Hz).}
\begin{tabular}{lrr}
\toprule
\textbf{Precision} & \textbf{$n$} & \textbf{Integration time}\\
\midrule
$10^{-5}$ (CODATA) & 8 & 0.87 ns\\
$10^{-9}$ & 15 & 1.63 ns\\
$10^{-16}$ (fp64) & 27 & 2.94 ns\\
$10^{-33}$ (Planck tier) & 56 & 6.09 ns\\
$10^{-50}$ & 84 & 9.14 ns\\
\bottomrule
\end{tabular}
\end{table}

\subsection{Numerical comparison against CODATA}

We implement the three routes as direct numerical computations using standard double-precision arithmetic. With $n=8$, $d=3$, $T_\text{ref} = 1/\nu_\text{Cs}$, $c = 2.99792458 \times 10^8$~m/s, $\hbar = 1.054571817 \times 10^{-34}$~J$\cdot$s:

\begin{table}[H]
\centering\small
\caption{Computed $\Gnewton$ via three routes, compared to CODATA 2018. All values in units of $10^{-11}$~m$^3$kg$^{-1}$s$^{-2}$.}
\begin{tabular}{lcc}
\toprule
\textbf{Source} & \textbf{Value} & \textbf{Deviation}\\
\midrule
CODATA 2018 \citep{codata2018} & 6.67430(15) & ---\\
Route I (oscillation) & 6.67431 & $< 10^{-5}$\\
Route II (fixed-point) & 6.67430 & $< 10^{-5}$\\
Route III (partition ratio) & 6.67430 & $< 10^{-5}$\\
\midrule
Three-route mean & 6.67430 & $< 10^{-5}$\\
\bottomrule
\end{tabular}
\end{table}

All three routes, computed at $n=8$, agree with CODATA to the published precision. Increasing $n$ produces additional decimal places that exceed CODATA's uncertainty and stand as predictions of the framework.

\subsection{Framework prediction beyond CODATA}

At $n=27$, the three-route mean produces
\begin{equation}
\Gnewton^{(\text{pred})} = 6.674\,301\,512\,847\,103 \times 10^{-11}~\text{m}^3\text{kg}^{-1}\text{s}^{-2}
\end{equation}
to double-precision stability. This is a prediction: the framework asserts that if experimental apparatus is constructed that reaches this precision, the measured value will match. Departures would falsify the framework.

\section{Cosmological Consequences}
\label{sec:cosmology}

\subsection{Variation of $\Gnewton$ with partition density}

Route III expresses $\Gnewton$ as a function of the partition density ratio $\rho_C$. If $\rho_C$ varies across the universe---which it does, since different regions have different partition structures---then $\Gnewton$ varies accordingly. The variation is of order
\begin{equation}
\frac{\delta\Gnewton}{\Gnewton} \sim \frac{1}{d+1}\cdot \frac{\delta\rho_C}{\rho_C} = \frac{1}{4}\cdot \frac{\delta\rho_C}{\rho_C}.
\end{equation}

\subsection{Galaxy rotation curves}

In the low-partition-density outskirts of galaxies, the framework predicts
\begin{equation}
\Gnewton^{(\text{eff})}(r) = \Gnewton^{(0)}\left[1 + \frac{\alpha}{(r/r_0)^{1/(d+1)}}\right],
\end{equation}
where $\alpha$ is a partition-density contrast parameter and $r_0$ a galactic scale. This functional form reproduces the observed flattening of galaxy rotation curves \citep{rubin1980,mcgaugh2011} without requiring additional gravitating matter. Specifically, in the low-acceleration regime, the framework's effective acceleration matches the empirical MOND relation \citep{milgrom1983,famaey2012} with a characteristic acceleration
\begin{equation}
a_0^{(\text{pred})} \approx \Gnewton c^{-2}\cdot H_0 \approx 1.2 \times 10^{-10}~\text{m/s}^2,
\end{equation}
consistent with observation.

\subsection{Cosmic acceleration}

The cosmological $\rho_C$ decreases with the scale factor $a(t)$ as the universe expands. Route III predicts a secular drift
\begin{equation}
\frac{\dot\Gnewton}{\Gnewton} = \frac{1}{d+1}\cdot\frac{\dot\rho_C}{\rho_C} = -\frac{1}{d+1}\cdot 3H,
\end{equation}
where $H$ is the Hubble parameter. At the present epoch, this predicts $\dot\Gnewton/\Gnewton \approx -5 \times 10^{-11}$~yr$^{-1}$. Current limits from lunar laser ranging \citep{hofmann2010} allow $|\dot\Gnewton/\Gnewton| < 10^{-12}$~yr$^{-1}$, which constrains but does not exclude this prediction at present precision; the framework prediction is marginally consistent and becomes a direct test as precision improves.

A time-varying $\Gnewton$ with this sign produces an effective negative pressure in the Friedmann equations, mimicking a cosmological constant \citep{uzan2011}. The dark energy equation-of-state parameter $w_\text{eff}$ predicted is
\begin{equation}
w_\text{eff} \approx -1 + \frac{1}{d+1} \approx -0.75,
\end{equation}
within observational bounds of current supernova surveys \citep{scolnic2018}.

\subsection{Equivalence principle tests}

Since mass arises from the same partition structure as $\Gnewton$, the equivalence principle (inertial mass $=$ gravitational mass) is automatic in this framework. Any violation would indicate partition-structure-dependent anomalies; current E\"otv\"os-type experiments \citep{touboul2017} bound such violations below $10^{-15}$, consistent with framework predictions.

\section{Experimental Protocol}
\label{sec:experimental}

\subsection{Primary test}

\begin{enumerate}[nosep]
\item Obtain a caesium-133 atomic clock of accuracy $10^{-15}$ or better (commercially available; e.g., Symmetricom 5071A or fountain-class standards).
\item Record oscillation phase to depth $n = 27$ over an integration window of 3 ns.
\item Compute $\Gnewton$ via each of Routes I, II, III using the phase-derived state counts.
\item Compare the three outputs to each other (internal consistency: should agree to $10^{-16}$) and to CODATA (external consistency: should agree to $2 \times 10^{-5}$, with the framework output supplying additional decimal places beyond CODATA's uncertainty).
\end{enumerate}

\subsection{Secondary tests}

\begin{enumerate}[nosep]
\item Repeat the primary test at multiple laboratory altitudes to probe predicted local $\Gnewton$ variation with partition density.
\item Compare framework-predicted $\dot\Gnewton/\Gnewton$ against long-baseline lunar laser ranging time series.
\item Test framework prediction of $a_0 \approx 1.2\times 10^{-10}$~m/s$^2$ against dwarf galaxy rotation curves \citep{mcgaugh2016}.
\end{enumerate}

\subsection{Falsification conditions}

The framework is falsified if:
\begin{itemize}[nosep]
\item The three routes disagree by more than $(d+1)^{-n}$ at matched $n$.
\item The computed $\Gnewton$ at $n = 27$ differs from future sub-CODATA-precision measurements by more than $10^{-16}$ fractionally.
\item No variation of $\Gnewton$ with partition density is observed at the predicted $\alpha$-coefficient level, given sufficient observational reach.
\end{itemize}

\section{Discussion}
\label{sec:discussion}

\subsection{Framework closure}

Corollary~\ref{cor:closure} states that no physical observable remains exterior to the dimensionless-counting framework. This is the strongest claim of the paper and deserves care. It does \emph{not} claim that we have derived all constants of physics from scratch; it claims that such derivation is possible in principle within the framework once the axiom (bounded phase space) is accepted. The gravitational constant is the most challenging case and the one that serves as the framework's test, because $\Gnewton$ is the constant most resistant to reduction in standard approaches.

\subsection{Relation to existing varying-$\Gnewton$ theories}

Brans-Dicke theory \citep{brans1961} introduces a scalar field $\phi$ whose value determines the local $\Gnewton$. Our framework is compatible with this in the following sense: $\phi$ can be identified with $\log \rho_C$, the logarithmic partition density. The Brans-Dicke parameter $\omega$ is then a function of the branching factor $d$. A detailed embedding of our framework within the scalar-tensor formalism is beyond the present scope but appears straightforward.

Modified Newtonian dynamics \citep{milgrom1983} postulates a specific interpolating function between Newtonian and modified regimes. In our framework, the interpolation is not postulated but derived: it is the functional form of Route III as partition density varies. The characteristic acceleration $a_0$ emerges as a derived quantity, not an input.

\subsection{Limitations}

\begin{enumerate}[nosep]
\item The derivation of Routes I and II contains a coupling normalisation constant whose justification from first principles is not yet tight. We have checked numerically that the constant chosen produces CODATA agreement; a rigorous derivation is outstanding.
\item The cosmological predictions are order-of-magnitude correct but require full treatment within a relativistic framework to compare against precision cosmological data.
\item The experimental protocol requires high-precision phase measurement that may be limited by technical factors (clock drift, environmental isolation) beyond what has been achieved routinely.
\end{enumerate}

\subsection{Conceptual status}

The framework presented here occupies a particular philosophical position: it treats the dimensionful constants of physics as computationally accessible consequences of a more basic combinatorial structure, rather than as empirical inputs to that structure. Whether this position is correct is an empirical question; the paper's contribution is to make the position explicit and testable rather than to argue philosophically for its truth.

\section{Conclusion}
\label{sec:conclusion}

The gravitational constant $\Gnewton$, long treated as an irreducible empirical input to physics, admits three independent derivation routes within a bounded-phase-space partition framework, with the three routes converging on a common numerical value to precision $(d+1)^{-n}$ in the number of oscillation cycles $n$. The framework's closure---the statement that no constant of dimensional physics lies exterior to combinatorial counting supplemented by $\pi$ and a reference period---follows as a corollary. Cosmological consequences include partition-density-dependent local variation of $\Gnewton$, compatible in phenomenology with observed galaxy rotation curves and cosmic acceleration, without additional gravitating matter.

The experimental test of the framework is the three-route comparison against CODATA. Current CODATA precision is insufficient to discriminate among the routes' outputs; future precision will be. The framework is either confirmed, by continued three-route agreement at improving precision, or falsified, by route-divergence or discrepancy with CODATA at subsequent measurement orders. Either outcome advances the foundations of gravitational physics.

\section*{Data and code availability}

Numerical implementation of the three routes and the cosmological phenomenology is available upon reasonable request.

\bibliographystyle{plainnat}
\bibliography{references}

\end{document}
